\section{Introduction}\label{sec:introduction}

The real embeddings of a totally real number field turn its ideals into
lattices. Multiplication by positive diagonal matrices changes the relative
sizes of the embedding coordinates while preserving covolume. The units
make each resulting diagonal orbit compact, and the ideal classes assemble
these orbits into a finite packet. The distribution of such packets connects
the arithmetic of number fields with the dynamics of higher-rank diagonal
groups. We prove equidistribution for totally real sextic fields having no
proper intermediate field.

\subsection{Packets and the main theorem}

Let
\[
 X_6=\SL_6(\Z)\backslash\SL_6(\R),\qquad
 A_6=\{\diag(e^{u_1},\ldots,e^{u_6}):\textstyle\sum_j u_j=0\}.
\]
We regard $X_6$ as the space of covolume-one lattices in $\R^6$, with
matrices acting on row vectors on the right. Let $m_6$ denote its
$\SL_6(\R)$-invariant probability. We normalize Haar measure on $A_6$
as $\dd u_1\cdots\dd u_5$, with $u_6=-\sum_{j=1}^5u_j$.

A degree-six number field $K$ is \emph{primitive} if there is no field
$F$ with $\Q\subsetneq F\subsetneq K$. Suppose $K$ is totally real,
choose an ordering $\sigma=(\sigma_1,\ldots,\sigma_6)$ of its real
embeddings, and write $\mathcal O_K$ for its ring of integers. For a
nonzero fractional $\mathcal O_K$-ideal $I$, set
\[
 \sigma(I)=\{(\sigma_1(d),\ldots,\sigma_6(d)):d\in I\},\qquad
 \Lambda_{I,\sigma}=\covol(\sigma(I))^{-1/6}\sigma(I).
\]
Choose one representative of each ordinary ideal class. Let
$\mathcal P_{K,\sigma}$ be the set of distinct compact orbits
\[
 \Lambda_{I,\sigma}wA_6,
 \qquad w=\diag(\epsilon_1,\ldots,\epsilon_6),\quad
 \epsilon_j\in\{1,-1\}.
\]
Here every sign matrix acts on the underlying lattice, including those
of determinant $-1$. Dirichlet's unit theorem gives compactness of the
orbits. Including all signs makes the packet independent of the chosen
ordinary ideal-class representatives.

For $\mathcal O\in\mathcal P_{K,\sigma}$, let $m_{\mathcal O}$ be
its invariant probability and $\vol_A(\mathcal O)$ its volume for the
specified Haar measure. The packet probability is
\begin{equation}\label{eq:packet-definition}
 \mu_{K,\sigma}=
 \frac{\displaystyle\sum_{\mathcal O\in\mathcal P_{K,\sigma}}
              \vol_A(\mathcal O)m_{\mathcal O}}
      {\displaystyle\sum_{\mathcal O\in\mathcal P_{K,\sigma}}
              \vol_A(\mathcal O)}.
\end{equation}

\begin{theorem}\label{thm:main}
Let $K_i$ be totally real primitive fields of degree six, with
$|\Disc(K_i)|\to\infty$, and choose any ordering $\sigma_i$ of the
embeddings of each field. Then
\[
 \int_{X_6} f\,\dd\mu_{K_i,\sigma_i}
       \longrightarrow\int_{X_6} f\,\dd m_6
       \qquad(f\in C_c(X_6)).
\]
There is no escape of mass: for every $\epsilon>0$ there is a compact
$C\subset X_6$ such that
$\mu_{K_i,\sigma_i}(C)\ge1-\epsilon$ for all sufficiently large $i$.
\end{theorem}

The theorem concerns each complete packet attached to a field. It allows
the fields and their embedding orderings to vary freely under the stated
hypotheses. The arithmetic role of primitivity is to make every
nonrational integer of $K_i$ generate the full field; this will yield
uniform separation in the dual lattice used to count thin tubes.

\subsection{History and the composite-degree obstruction}

For real quadratic fields, the corresponding packets project to packets
of closed geodesics on the modular surface. Linnik's ergodic method,
including Skubenko's treatment of the one-sheeted hyperboloid, gave
equidistribution under an auxiliary splitting condition
\cite[Chapter~VI]{Linnik1968}; see also the historical statements in
\cite[Theorems~1.1--1.2]{ELMV2012}. Duke removed this condition for
fundamental discriminants by analytic methods
\cite[Theorem~1]{Duke1988}. Einsiedler, Lindenstrauss, Michel, and
Venkatesh later gave a dynamical proof formulated for all positive
nonsquare order discriminants~\cite[Theorem~2.3]{ELMV2012}.
In higher degree the compact diagonal orbits have dimension greater
than one, and the dynamics of the full diagonal group provides
additional rigidity.

Einsiedler, Katok, and Lindenstrauss proved that, on
$\SL_d(\Z)\backslash\SL_d(\R)$ for $d\ge3$, an ergodic probability
invariant under the full positive diagonal group and having positive
entropy for a diagonal flow is algebraic~\cite[Theorem~1.3]{EKL2006}.
Einsiedler, Lindenstrauss, Michel,
and Venkatesh developed the connection between arithmetic separation of
periodic torus orbits and entropy of their limiting measures
\cite{ELMV2009}. Their subsequent cubic equidistribution theorem
combined measure classification, subconvexity, and local harmonic
analysis~\cite{ELMV2011}. These works distinguish two tasks that remain
central here: establishing nonescape and entropy, and identifying the
algebraic components of a limit.

The second task is particularly important in composite degree
\cite[Section~1.6.3]{ELMV2011}. In prime
degree, the positive-entropy classification gives Haar measure directly
\cite[Corollary~1.4]{EKL2006}. In degree six there are proper block
subgroups containing the full diagonal group whose closed orbits can
carry finite invariant measures. Primitivity of the fields in the
sequence does not, by itself, rule out concentration near such orbits.
The proof must retain enough information from the sequence to exclude
these limiting components.

Subsequent work has sharpened the entropy estimates in higher degree.
Khayutin obtained bounds for maximal-order packets whose Galois closures
act two-transitively on the embeddings~\cite[Theorem~1.1]{Khayutin2019}.
Wieser and Yang obtained entropy bounds for quartic packets of maximal
type with quadratic subfields, under a fixed split-place condition and
additional discriminant-growth hypotheses~\cite{WieserYang2026}.
These results clarify the arithmetic information available in composite
degree; identifying every component of a limit remains a separate step
in the argument here.

The arithmetic preparation adapts the ordinary-vector method used for
prime-degree packets~\cite[Sections~2--6]{PrimeDegreePackets} and primitive
quartic packets~\cite[Section~3]{QuarticPackets}. We give the needed estimates in full.
They use Stark's effective zero-free information
\cite{Stark1974} and the short-interval bounds of Shiu and Pollack
\cite{Shiu1980,Pollack2020}. The quartic work treats its additional
obstruction using exterior-square arithmetic. Here an ordinary vector
and a dual vector produce a codifferent element, and primitivity gives
the separation needed to estimate its partial traces.

\subsection{The intermediate-resolution argument}

The main additional argument concerns conditional measures that retain
information at growing, but separated, resolutions. We describe its
purpose before introducing the construction.

Let $\mu_i=\mu_{K_i,\sigma_i}$. Ordinary-vector estimates prove
tightness and give every probability limit $\mu$ an $O(r^6)$ bound
for the mass of radius-$r$ balls, uniformly when their centers range
over a compact set. Since $A_6$ has dimension five, this bound
forces positive entropy in almost every $A_6$-ergodic component.
Theorem~\ref{thm:ekl} then makes those components homogeneous.
Every proper component is supported on points fixed by a nonidentity
diagonal matrix with algebraic entries, so only a countable family of
fixed-point sets must be excluded.

The arithmetic estimate used for this exclusion concerns two-sided
tubes, with $O$ a sufficiently small identity neighborhood,
\[
 x\bigl(a(tb)Oa(-tb)\cap a(-tb)Oa(tb)\bigr),
 \qquad a(b)=\diag(e^{b_1},\ldots,e^{b_6}),\quad\sum_jb_j=0.
\]
If the coordinates of $b$ are separated across a prescribed cut of the
indices, ordinary and dual vectors bound the packet mass of such a
tube by $e^{-\eta t}$ for a fixed $\eta>0$. The estimate holds throughout a growing interval
$[T_{-,i},T_{+,i}]$ with $T_{+,i}/T_{-,i}\to\infty$.
The breadth of this interval is used throughout the proof.

For $j\ne k$, let $U_{jk}=\{1+sE_{jk}:s\in\R\}$.
Exact conditional measures of a finite packet along these root groups
are supported at a single point. They therefore cannot express the
spread detected by the arithmetic tube bound. Instead, in a local box
$(q,h)$ for a unipotent group, we condition the plaque coordinate $h$
on a small cell of the transverse coordinate $q$. We choose logarithmic
cell depths inside the growing interval so that the conditional
probabilities change negligibly across a much wider internal range.
In their joint limits with the sampled lattice, the plaque point still
has the recorded probability as its conditional law, given the transverse
coordinate and that one probability. This sampling identity will detect
atoms on sets meeting each root plaque in only countably many points.
The limiting probabilities also retain their behavior under suitable
long diagonal motions.

Product structure for genuine leafwise measures and the use of
noncommuting root foliations are basic ingredients of the higher-rank
entropy method of Einsiedler and Katok
\cite[Proposition~8.3]{EK2003}~\cite[Theorems~8.4--8.5]{EK2005};
see also~\cite[Corollary~8.8 and Theorem~9.8]{EL2010}. The point here is
to prove the corresponding structure for the conditional probabilities
formed before taking the limit. If a diagonal element fixes one root
group $U$ and contracts a complementary unipotent subgroup $W$ generated
by root groups and normal in $WU$,
the limiting measure on $WU$ is the product of the limiting measures
on $W$ and $U$. Contraction identifies the $U$ factor. A separate
conditional-information argument shows that adding a fine $U$-cell label
barely changes the law of any fixed further subdivision of a typical
$W$ tile. Expanding that tile turns this comparison into one for bounded
$W$ tests and identifies the $W$ factor. Both arguments are needed, and neither
follows by simply disintegrating the weak limit.

Call $j\to k$ active when the resulting $U_{jk}$ measure has support
beyond the identity. The product rule makes activity transitive under
composition of roots. It also permits an entropy argument: if crossing
roots are inactive along a positive proportion of long trajectories,
then their names can be encoded using only the slowly expanding
directions inside the two sides of a cut. This contradicts the tube
bound. Long motions in the diagonal wall $b_j=b_k$ fixing $U_{jk}$ leave
its activity unchanged. Combining the resulting one-dimensional functions
of the differences $b_k-b_j$ upgrades activity along trajectories to full
activity at almost every limiting state. A proper homogeneous
fixed-point piece forces an inactive root and is therefore excluded.

The construction separates the arithmetic estimate from this last
argument. Its potentially reusable ingredients are the comparison of
conditional probabilities across a broad range of transverse depths,
their product rule, and the passage from entropy bounds for long segments
starting in any fixed positive proportion of a packet to full root activity.

Section~\ref{sec:arithmetic} proves the arithmetic bounds, and
Section~\ref{sec:homogeneous} identifies the homogeneous alternatives.
Sections~\ref{sec:conditioning} and~\ref{sec:product} construct and
factor the intermediate conditional measures.
Section~\ref{sec:activity} records their root-group consequences.
Sections~\ref{sec:segment-entropy} and~\ref{sec:combination} prove the
long-segment and full-activity assertions, respectively, and
Section~\ref{sec:assembly} concludes the proof.

\subsection{Conventions}

For the remainder of the paper write $n=6$,
$G=\SL_n(\R)$, $\Gamma=\SL_n(\Z)$, $X=\Gamma\backslash G$,
and $A=A_n$. Let $L=\{b\in\R^n:\sum_j b_j=0\}$, identified with the
traceless diagonal matrices, and
$a(b)=\diag(e^{b_1},\ldots,e^{b_n})$. We use the right action
throughout. Thus motion by $a(b)$ changes a relative displacement $h$
to $a(-b)ha(b)$.
Matrix neighborhoods of the identity and smooth metrics are interchangeable
on fixed compact sets. When a metric on all of $X$ is needed, we use
the quotient of a left-invariant Riemannian metric on $G$.
The notation $B(r)$ denotes a sufficiently small identity neighborhood
of radius $r$. Implicit constants may depend on fixed compact sets,
charts, and $n$, but not on the varying field. All passages to
subsequences preserve the probability limit under consideration.
